% problem-000118 7 配位化学与氧化态
\section*{7. 配位化学与氧化态}
菲啶(Phend)形成的铜配合物因在太阳能转化、发光⼆极管、荧光探针和抗癌荧光等光伏领域的应⽤⽽受到研究⼈员的重视。某实验将 0.089 g 菲啶(Phend)和 0.131 g $\mathrm { P P h } _ { 3 }$ 的25 mL甲醇溶液室温下搅拌10分钟，然后加⼊含0.049g氯化亚铜的25 mL甲醇溶液，室温搅拌3⼩时得到⻩绿⾊沉淀。过滤，滤液室温下放置10天，随着溶剂的缓慢蒸发，得到0.253 g⻩⾊晶体 $\mathbf { A } _ { \circ }$ 。晶体 A 的元素分析结果为(%)：C 68.74 (68.89), H 4.53(4.47), N 2.27 (2.60)，括号内为理论值。晶体结构表明Cu的配位数为4。菲啶的结构⻅图5。

（图：）
图5 菲啶结构

\subsection*{7-1}
请通过计算和分析，写出⻩⾊晶体A的化学式。

\subsection*{7-2}
画出A的结构，写出形成体的价电⼦构型和杂化⽅式。

\subsection*{7-3}
画出A可能的异构体结构。
